% Source: source/普通高考/2017/2017全国1理(河南,河北,山西,江西,湖北,湖南,广东,安徽,福建).pdf, page 1.
% The source flowchart is vector line art; pdfimages -list found no embedded image objects.
% Original-side comparison crop: tmp/redraw_flowchart_2017_national_paper_1_science/original_q8_fig1.png,
% cropped from the ungridded page after 100px coarse and 50px fine grid inspection.
% Elements: rounded start/end boxes, input/output parallelograms, assignment rectangle,
% blank assignment rectangle, decision diamond, directed connectors, and 是/否 branch labels.
% Node -> directed-edge list from the source: 开始 -> 输入 n=0 -> A=3^n-2^n -> decision;
% decision --是--> blank assignment -> A=3^n-2^n; decision --否--> 输出 n -> 结束.
% This is a schematic flowchart.
% solid segments: all box borders, the central arrows, and the loop connectors.
% dashed segments: none.
% labels: math labels centered in boxes; 是 sits below the left decision exit and 否 to the
% right of the downward exit; the blank boxes intentionally contain no added text.
\begin{tikzpicture}[
    x=0.8cm,
    y=0.8cm,
    >=stealth,
    line width=0.45pt,
    every node/.style={inner sep=1pt, align=center, anchor=center,
        execute at begin node=\setlength{\parindent}{0pt}},
]
    % Main vertical algorithm path.
    \node[draw, rounded corners=6pt, minimum width=1.00cm, minimum height=0.56cm]
        (start) at (0,8.20) {开始};
    \node[draw, trapezium, trapezium left angle=70, trapezium right angle=110,
        minimum width=1.20cm, minimum height=0.62cm]
        (input) at (0,6.92) {输入 \(n=0\)};
    \node[draw, minimum width=2.20cm, minimum height=0.55cm]
        (assign) at (0,5.40) {\(A=3^n-2^n\)};
    \node[draw, diamond, aspect=4.0, minimum width=3.40cm, minimum height=0.85cm]
        (test) at (0,4.10) {};
    \node[draw, trapezium, trapezium left angle=70, trapezium right angle=110,
        minimum width=0.95cm, minimum height=0.60cm]
        (output) at (0,2.54) {输出 \(n\)};
    \node[draw, rounded corners=6pt, minimum width=1.00cm, minimum height=0.56cm]
        (finish) at (0,1.18) {结束};

    % Empty assignment box on the loop branch.
    \node[draw, minimum width=1.70cm, minimum height=0.58cm]
        (loopassign) at (-3.25,4.98) {};

    % Downward path and output branch.
    \draw[->] (start.south) -- (input.north);
    % Arrowheads stop in the clear corridor; final shafts enter the boxes.
    \draw[->] (input.south) -- (0,6.15);
    \draw[-] (0,6.15) -- (assign.north);
    \draw[->] (assign.south) -- (test.north);
    \draw[->] (test.south) -- (output.north)
        node[pos=0.45, right=2pt] {否};
    \draw[->] (output.south) -- (finish.north);

    % 是 branch: left tip -> empty assignment -> return to A.
    \draw (test.west) -- (-3.25,4.10) -- (-3.25,4.62);
    \draw[->] (-3.25,4.62) -- (loopassign.south);
    \node[inner sep=0pt] at (-2.10,3.72) {是};
    \draw[->] (loopassign.north) -- (-3.25,6.15) -- (-1.50,6.15);
    \draw[-] (-1.50,6.15) -- (assign.north);
\end{tikzpicture}
